\documentclass[../../main.tex]{subfiles}
\begin{document}
\section*{习题}\label{section:ch11-习题}

\begin{exercise}\label{exercise:ch11-01}
设 $V$ 与 $W$ 为有限维向量空间，$A\colon V\to W$ 为任意线性映射。
\begin{enumerate}[(a)]
\item 证明下图交换：
\begin{align*}
\begin{array}{ccc}
V & \xrightarrow{A} & W\\
\zeta_V\downarrow && \downarrow\zeta_W\\
V^{**} & \xrightarrow{(A^*)^*} & W^{**},
\end{array}
\end{align*}
其中 $\zeta_V$ 与 $\zeta_W$ 分别是 $V$、$W$ 的典范同构。
\item 证明不可能给每个有限维向量空间 $V$ 指定一个同构 $\beta_V\colon V\to V^*$，使得对每个线性映射 $A\colon V\to W$，下图交换：
\begin{align*}
\begin{array}{ccc}
V & \xrightarrow{A} & W\\
\beta_V\downarrow && \downarrow\beta_W\\
V^* & \xrightarrow{A^*} & W^*.
\end{array}
\end{align*}
\end{enumerate}
\end{exercise}

\begin{exercise}\label{exercise:ch11-02}
设 $V$ 为无限维实向量空间，$V^*$ 表示从 $V$ 到 $\mathbb{R}$ 的所有线性泛函构成的向量空间。
\begin{enumerate}[(a)]
\item 证明存在从 $V$ 到 $V^*$ 的单线性映射。（提示：见习题~B.5。）
\item 证明不存在从 $V^*$ 到 $V$ 的单线性映射。（提示：若 $S$ 是 $V$ 的一组基，证明 $V^*$ 中存在一个与 $S$ 的幂集等势的线性无关子集。）
\item 用 (a) 与 (b) 证明 $V$ 与 $V^{**}$ 之间不存在同构。
\end{enumerate}
（用于第~489~页。）
\end{exercise}

\begin{exercise}\label{exercise:ch11-03}
设 $\mathsf{Vec}_1$ 为习题~10-8~中的有限维向量空间与线性同构范畴。定义函子 $F\colon\mathsf{Vec}_1\to\mathsf{Vec}_1$：对向量空间 $V$ 令 $F(V)=V^*$，对同构 $A$ 令 $F(A)=(A^{-1})^*$。证明 $F$ 是光滑共变函子，且对每个 $M$，$F(TM)$ 与 $T^*M$ 典范光滑同构。
\end{exercise}

\begin{exercise}\label{exercise:ch11-04}
设 $M$ 为带边或不带边的光滑流形，$p\in M$。令 $I_p$ 为 $C^\infty(M)$ 中在 $p$ 处取零的光滑函数构成的子空间，$I_p^2$ 为由形如 $fg$（$f,g\in I_p$）的函数张成的 $I_p$ 的子空间。
\begin{enumerate}[(a)]
\item 证明 $f\in I_p^2$ 当且仅当在任意光滑局部坐标中，$f$ 在 $p$ 处的一阶 Taylor 多项式为零。（正因如此，$I_p^2$ 中的函数称为\textbf{二阶消失}。）
\item 定义映射 $\Phi\colon I_p\to T_p^*M$ 为 $\Phi(f)=df_p$。证明 $\Phi$ 在 $I_p^2$ 上的限制为零，且 $\Phi$ 下降为从 $I_p/I_p^2$ 到 $T_p^*M$ 的向量空间同构。
\end{enumerate}
（注：习题~3-8~表明 $p$ 处的切向量可视为光滑曲线的等价类，即从 $\mathbb{R}$（的子集）到 $M$ 的光滑映射。本习题表明 $p$ 处的余向量可对偶地视为从 $M$ 到 $\mathbb{R}$ 的光滑函数的等价类：余向量是在 $p$ 处消失的光滑函数的等价类，两个这样的函数若相差一个二阶消失的函数则视为等价。在光滑流形理论的一些处理中，$T_p^*M$ 首先这样定义，再把 $T_pM$ 定义为对偶空间 $(I_p/I_p^2)^*$。）
\end{exercise}

\begin{exercise}\label{exercise:ch11-05}
对任意光滑流形 $M$，证明 $T^*M$ 是平凡向量丛当且仅当 $TM$ 是平凡的。
\end{exercise}

\begin{exercise}\label{exercise:ch11-06}
设 $M$ 为光滑 $n$ 流形，$p\in M$，$y^1,\dots,y^k$ 为定义在 $M$ 中 $p$ 某邻域上的光滑实值函数。证明下列命题：
\begin{enumerate}[(a)]
\item 若 $k=n$ 且 $(dy^1|_p,\dots,dy^n|_p)$ 是 $T_p^*M$ 的一组基，则 $(y^1,\dots,y^n)$ 是 $M$ 在 $p$ 某邻域上的光滑坐标。
\item 若 $(dy^1|_p,\dots,dy^k|_p)$ 是线性无关的 $k$ 元余向量组且 $k<n$，则存在光滑函数 $y^{k+1},\dots,y^n$，使 $(y^1,\dots,y^n)$ 是 $M$ 在 $p$ 某邻域上的光滑坐标。
\item 若 $(dy^1|_p,\dots,dy^k|_p)$ 张成 $T_p^*M$，则存在指标 $i_1,\dots,i_n$，使 $(y^{i_1},\dots,y^{i_n})$ 是 $M$ 在 $p$ 某邻域上的光滑坐标。
\end{enumerate}
（用于第~584~页。）
\end{exercise}

\begin{exercise}\label{exercise:ch11-07}
下面各题中，$M$ 与 $N$ 为光滑流形，$F\colon M\to N$ 为光滑映射，$\omega\in\mathfrak X^*(N)$。计算每种情形的 $F^*\omega$。
\begin{enumerate}[(a)]
\item $M=N=\mathbb{R}^2$，$F(s,t)=(st,e^t)$，$\omega=x\,dy-y\,dx$。
\item $M=\mathbb{R}^2$，$N=\mathbb{R}^3$，$F(\varphi,\psi)=((\cos\psi+2)\cos\varphi,(\cos\psi+2)\sin\varphi,\sin\psi)$，$\omega=z^2 dx$。
\item $M=\{(s,t)\in\mathbb{R}^2:s^2+t^2<1\}$，$N=\mathbb{R}^3\smallsetminus\{0\}$，$F(s,t)=(s,t,\sqrt{1-s^2-t^2})$，$\omega=(1-x^2-y^2)dz$。
\end{enumerate}
\end{exercise}

\begin{exercise}\label{exercise:ch11-08}
(a) 设 $F\colon M\to N$ 是微分同胚，令 $dF^*\colon T^*N\to T^*M$ 为在每个余切空间 $T_q^*N$ 上等于 $dF^*_{F^{-1}(q)}$ 的映射。证明 $dF^*$ 是光滑丛同态。

(b) 设 $\mathsf{Diff}_1$ 为以光滑流形为对象、仅以微分同胚为态射的范畴；$\mathsf{VB}$ 为以光滑向量丛为对象、光滑丛同态为态射的范畴。证明指派 $M\mapsto T^*M$、$F\mapsto dF^*$ 定义了从 $\mathsf{Diff}_1$ 到 $\mathsf{VB}$ 的反变函子，称为\textbf{余切函子}。（用于第~592~页。）
\end{exercise}

\begin{exercise}\label{exercise:ch11-09}
设 $f\colon\mathbb{R}^3\to\mathbb{R}$ 为 $f(x,y,z)=x^2+y^2+z^2$，$F\colon\mathbb{R}^2\to\mathbb{R}^3$ 为球极投影的逆映射：
\begin{align*}
F(u,v)=\left(\frac{2u}{u^2+v^2+1},\frac{2v}{u^2+v^2+1},\frac{u^2+v^2-1}{u^2+v^2+1}\right).
\end{align*}
分别计算 $F^*df$ 与 $d(f\circ F)$，并验证二者相等。
\end{exercise}

\begin{exercise}\label{exercise:ch11-10}
下列各情形中，$M$ 为光滑流形，$f\colon M\to\mathbb{R}$ 为光滑函数。计算 $df$ 的坐标表示，并确定使 $df_p=0$ 的所有点 $p\in M$。
\begin{enumerate}[(a)]
\item $M=\{(x,y)\in\mathbb{R}^2:x>0\}$，$f(x,y)=x/\sqrt{x^2+y^2}$，用标准坐标 $(x,y)$。
\item $M$ 与 $f$ 如 (a)，此次用极坐标 $(r,\theta)$。
\item $M=\mathbb{S}^2\subseteq\mathbb{R}^3$，$f(p)=z(p)$（$p$ 作为 $\mathbb{R}^3$ 中点的 $z$ 坐标），用南、北极球极坐标（习题~1-7）。
\item $M=\mathbb{R}^n$，$f(x)=|x|^2$，用标准坐标。
\end{enumerate}
\end{exercise}

\begin{exercise}\label{exercise:ch11-11}
设 $M$ 为光滑流形，$C\subseteq M$ 为嵌入子流形。设 $f\in C^\infty(M)$，$p\in C$ 是 $f$ 在 $C$ 上取局部极大或极小的点。给定 $C$ 在 $p$ 某邻域 $U$ 上的光滑定义函数 $\Phi\colon U\to\mathbb{R}^k$，证明存在实数 $\lambda_1,\dots,\lambda_k$（称为\textbf{Lagrange 乘子}）使得
\begin{align*}
df_p=\lambda_1\,d\Phi_1|_p+\cdots+\lambda_k\,d\Phi_k|_p.
\end{align*}
\end{exercise}

\begin{exercise}\label{exercise:ch11-12}
证明连通光滑流形中任意两点都可用光滑曲线段连接。
\end{exercise}

\begin{exercise}\label{exercise:ch11-13}
光滑曲线段 $\gamma\colon[a,b]\to\mathbb{R}^n$ 的长度定义为通常积分的值
\begin{align*}
L(\gamma)=\int_a^b|\gamma'(t)|\,dt.
\end{align*}
证明不存在光滑余向量场 $\omega\in\mathfrak X^*(\mathbb{R}^n)$，使对每条光滑曲线 $\gamma$ 都有 $\int_\gamma\omega=L(\gamma)$。
\end{exercise}

\begin{exercise}\label{exercise:ch11-14}
考虑 $\mathbb{R}^3$ 上的两个余向量场：
\begin{align*}
\omega&=-\frac{4z\,dx}{(x^2+1)^2}+\frac{2y\,dy}{y^2+1}+\frac{2x\,dz}{x^2+1},\\
\eta&=-\frac{4xz\,dx}{(x^2+1)^2}+\frac{2y\,dy}{y^2+1}+\frac{2\,dz}{x^2+1}.
\end{align*}
\begin{enumerate}[(a)]
\item 写出并计算每个余向量场沿从 $(0,0,0)$ 到 $(1,1,1)$ 的直线段的线积分。
\item 判断这两个余向量场中哪一个是恰当的。
\item 对每个恰当的，求一个势函数并用它重算线积分。
\end{enumerate}
\end{exercise}

\begin{exercise}[向量场的线积分]\label{exercise:ch11-15}
设 $X$ 为开集 $U\subseteq\mathbb{R}^n$ 上的光滑向量场。给定分段光滑曲线段 $\gamma\colon[a,b]\to U$，定义 $X$ 沿 $\gamma$ 的线积分 $\int_\gamma X\cdot ds$ 为
\begin{align*}
\int_\gamma X\cdot ds=\int_a^b X_{\gamma(t)}\cdot\gamma'(t)\,dt,
\end{align*}
右端点积表示在 $\gamma(t)$ 处切向量（等同于 $\mathbb{R}^n$ 中元素）之间的 Euclid 点积。一个\textbf{保守向量场}是绕每条分段光滑闭曲线的线积分为零的向量场。
\begin{enumerate}[(a)]
\item 证明 $X$ 保守当且仅当存在光滑函数 $f\in C^\infty(U)$ 使 $X=\mathrm{grad}\,f$。（提示：考虑由 $\omega_x(v)=X_x\cdot v$ 定义的余向量场 $\omega$。）
\item 设 $n=3$。证明若 $X$ 保守，则 $\mathrm{curl}\,X=0$，其中
\begin{align*}
\mathrm{curl}\,X=\left(\frac{\partial X^3}{\partial x^2}-\frac{\partial X^2}{\partial x^3}\right)\frac{\partial}{\partial x^1}
+\left(\frac{\partial X^1}{\partial x^3}-\frac{\partial X^3}{\partial x^1}\right)\frac{\partial}{\partial x^2}
+\left(\frac{\partial X^2}{\partial x^1}-\frac{\partial X^1}{\partial x^2}\right)\frac{\partial}{\partial x^3}.
\end{align*}
\item 证明若 $U\subseteq\mathbb{R}^3$ 是星形的，则 $X$ 在 $U$ 上保守当且仅当 $\mathrm{curl}\,X=0$。
\end{enumerate}
\end{exercise}

\begin{exercise}\label{exercise:ch11-16}
设 $M$ 为正维数的紧流形。证明 $M$ 上每个恰当余向量场在 $M$ 的每个连通分支中至少在两点处为零。
\end{exercise}

\begin{exercise}\label{exercise:ch11-17}
设 $\mathbb{T}^n=\mathbb{S}^1\times\cdots\times\mathbb{S}^1\subseteq\mathbb{C}^n$ 为 $n$ 环面。对每个 $j=1,\dots,n$，令 $\gamma_j\colon[0,1]\to\mathbb{T}^n$ 为曲线段
\begin{align*}
\gamma_j(t)=(1,\dots,e^{2\pi\mathrm{i}t},\dots,1)\quad\text{（$e^{2\pi\mathrm{i}t}$ 在第 $j$ 位）}.
\end{align*}
证明 $\mathbb{T}^n$ 上的闭余向量场 $\omega$ 恰当当且仅当对 $j=1,\dots,n$ 有 $\int_{\gamma_j}\omega=0$。（提示：先考虑 $\varepsilon_n^*\omega$，其中 $\varepsilon_n\colon\mathbb{R}^n\to\mathbb{T}^n$ 是光滑覆叠映射 $\varepsilon_n(x^1,\dots,x^n)=(e^{2\pi\mathrm{i}x^1},\dots,e^{2\pi\mathrm{i}x^n})$。）
\end{exercise}
\end{document}
